\section{Finite support calculus}
\label{sec:finite-support}

Fix an atom carrier $A$ and selected actions of $\Perm(A)$ on $X$ and $Y$.
The atom action is canonical, $\pi\act a=\pi(a)$, and the universes of
all three carriers are independent. Finite support is a property of an
individual element; it does not require every element of its carrier to
have finite support. Infinitude enters the calculus only at intersection.

\subsection{Supporting bounds}
\label{sec:finite-support-definition}

\begin{definition}
A finite set $S\subseteq A$ \emph{supports} $x\in X$ when
\[
  \operatorname{Supports}(S,x)
  \quad\Longleftrightarrow\quad
  \forall\pi\in\Perm(A),\quad
  \bigl(\forall a\in S,\ \pi(a)=a\bigr)
  \Longrightarrow \pi\act x=x.
\]
The element $x$ is \emph{finitely supported} if such a finite $S$ exists.
\end{definition}

The definition makes a supporting set a sufficient bound, without choosing
one or asserting minimality. It requires pointwise fixation of $S$;
setwise preservation may permute its members. Nor does fixation of $x$
imply pointwise fixation of a supporting set. That reverse implication
belongs to the stronger notion of strong support. Support also depends on
the action: for example, support on bare $\Perm(A)$ refers to left
multiplication, not silently to conjugation
(Example~\ref{ex:group-actions}).

The Lean predicate \texttt{Supports S x}, with \texttt{S : Finset A},
is an abbreviation for Mathlib's two-carrier predicate
\begin{verbatim}
  MulAction.Supports (Perm A) (S : Set A) x
\end{verbatim}
This reuses its pointwise-fixer semantics with the canonical action on
atoms and the selected action on values. The atom sort is determined by
$S$, whereas \texttt{FinitelySupported A x} takes it explicitly.
Both predicates inhabit \texttt{Prop}, so neither supplies an observable
choice of bound. Neither definition requires decidable equality
or infinitude of $A$.

\subsection{Elementary bounds and equivariant images}
\label{sec:finite-support-basic}
\label{sec:finite-support-products}

\begin{proposition}[Elementary support bounds]
\label{prop:support-basic}
\label{prop:support-products}
\label{prop:finitely-supported-basic}
For finite bounds $S,T\subseteq A$:
\begin{enumerate}
\item If $S$ supports $x$ and $S\subseteq T$, then $T$ supports $x$.
\item The empty set supports $x$ exactly when every permutation fixes $x$.
\item If $f:X\to Y$ is equivariant and $S$ supports $x$, then $S$ supports
  $f(x)$.
\item Under the componentwise action, $S$ supports $(x,y)$ exactly when it
  supports both components. If $S$ supports $x$ and $T$ supports $y$,
  then $S\cup T$ supports $(x,y)$.
\end{enumerate}
Consequently, equivariant images and products preserve finite supportedness.
None of these conclusions requires infinitude.
\end{proposition}

\begin{proof}
A permutation fixing $T$ pointwise also fixes $S\subseteq T$; for the
empty set the pointwise condition is vacuous. Equivariance and support give
$\pi\act f(x)=f(\pi\act x)=f(x)$.
For products, fixation is componentwise; enlarge separate bounds to
$S\cup T$. Existential witnesses for the image and product are the same
bound and the union, respectively.
\end{proof}

The image theorem is a bound, not an equality of least supports. For
example, a constant map into $\Discrete_A(\mathbb B)$ is equivariant.
Its value inherits the bound $\{a\}$ of an atom $a$, but also has the
strictly smaller empty bound. The image theorem needs only the
ordinary-function predicate $\Equivariant_A(f)$, without an action or a
support certificate on the function itself.

\subsection{Support transport by conjugation}
\label{sec:support-transport}

\begin{proposition}[Conjugation transport]
\label{prop:support-transport}
For every $\pi\in\Perm(A)$ and finite $S\subseteq A$,
\[
  \operatorname{Supports}(\pi[S],\pi\act x)
  \quad\Longleftrightarrow\quad
  \operatorname{Supports}(S,x).
\]
No infinitude or commuting-action hypothesis is required.
\end{proposition}

\begin{proof}
Suppose $S$ supports $x$ and $\sigma$ fixes $\pi[S]$ pointwise.
For $a\in S$, one has $\sigma(\pi(a))=\pi(a)$, so
$(\pi^{-1}\sigma\pi)(a)=a$.
Thus $\pi^{-1}\sigma\pi$ fixes $S$ pointwise and hence fixes $x$.
The action laws give
\[
  \sigma\act(\pi\act x)
  =\pi\act\bigl((\pi^{-1}\sigma\pi)\act x\bigr)
  =\pi\act x.
\]
For the converse, apply this implication to $\pi^{-1}$, the bound
$\pi[S]$ and the element $\pi\act x$, then cancel the inverse actions.
\end{proof}

The formal statement is \path{supports_smul_iff}. Conjugation is essential
to this proof: renaming transports the pointwise fixer of $S$ to that of
$\pi[S]$. Mathlib's \path{MulAction.Supports.smul} instead assumes
commuting scalar actions on both carriers, a condition that general
permutations do not satisfy. The conjugation argument accommodates their
noncommutativity directly.

\begin{proposition}[Existential transport]
\label{prop:finitely-supported-transport}
For every $\pi\in\Perm(A)$,
\[
  \operatorname{FinitelySupported}_A(\pi\act x)
  \quad\Longleftrightarrow\quad
  \operatorname{FinitelySupported}_A(x).
\]
\end{proposition}

Transport sends a witness $S$ to $\pi[S]$; applying $\pi^{-1}$ gives the
converse. Thus finite supportedness is
invariant under renaming even when no bound is selected as data.

\subsection{Swaps and intersection}
\label{sec:support-swap-criterion}
\label{sec:support-intersection}

\begin{proposition}[Support tested by swaps]
\label{prop:support-swap-criterion}
A finite $S\subseteq A$ supports $x$ if and only if
\[
  \forall a,b\in A,\quad a\notin S\ \land\ b\notin S
  \Longrightarrow\swap a b\act x=x.
\]
Equal endpoints are allowed, and no infinitude premise is required.
\end{proposition}

This specializes Proposition~\ref{prop:swap-invariance} to a finite bound.
It turns a quantification
over all finite permutations into a local condition on two atoms.
Intersection then reduces to finding a spare atom joining two such local
conditions.

\begin{proposition}[Intersection of finite support bounds]
\label{prop:finite-support-intersection}
Assume $A$ is infinite. If finite $S,T\subseteq A$ both support $x$,
then $S\cap T$ supports $x$.
\end{proposition}

\begin{proof}
By the swap criterion, it suffices to show that $\swap a b$ fixes $x$
whenever $a,b\notin S\cap T$. The case $a=b$ is immediate. Otherwise,
choose
\[
  c\notin S\cup T\cup\{a,b\}.
\]
Such an atom exists because the excluded set is finite. The atom $a$ lies
outside at least one of $S,T$, and $c$ lies outside both; therefore
$\swap a c$ fixes $x$. Similarly, $\swap c b$ fixes $x$. Conjugation gives
\[
  \swap a c\swap c b\swap a c
  =\swap{\swap a c(c)}{\swap a c(b)}
  =\swap a b.
\]
All three factors on the left fix $x$, so their product does too.
\end{proof}

In \path{supports_inter}, infinitude is used only to obtain $c$ by
Mathlib's \path{Finset.exists_notMem}. There is no assumption that all
values of $X$ are finitely supported. The role of the atom hypothesis is
mathematical, as the following example shows.

\begin{example}[Intersection can fail on finite atoms]
\label{ex:finite-atom-intersection}
Let $A=\mathbb B$ carry its canonical action, and take
$x=\mathrm{false}$. Both $\{\mathrm{false}\}$ and
$\{\mathrm{true}\}$ support $x$, but their empty intersection does not.
\end{example}

\begin{proof}
The first bound is immediate. A permutation fixing $\mathrm{true}$ also
fixes $\mathrm{false}$: injectivity prevents it from sending both atoms to
$\mathrm{true}$, and these are the only atoms. This gives the second
bound. Yet $\swap{\mathrm{false}}{\mathrm{true}}$ moves $x$, so the empty
set cannot support it.
\end{proof}

The example uses the same finite-permutation group and support definition
as the infinite-atom theorem. It shows why a sufficient support bound
cannot be identified with a least bound and why intersection needs an
atom hypothesis.

\subsection{Assumptions and support witnesses}
\label{sec:support-assumptions}

The mathematical distinction between a supplied finite bound and existence
of a bound has a useful consequence for the Lean interface. Visible union,
image and intersection expressions require decidable equality. When such
a set is constructed only as an existential witness, classical equality
can remain local to the proof: finite supportedness still carries no
computational choice of bound. The assumptions therefore separate as follows.

\begin{center}
\small
\begin{tabular}{@{}>{\raggedright\arraybackslash}p{0.66\linewidth}
  >{\raggedright\arraybackslash}p{0.28\linewidth}@{}}
\hline
Statements & Atom assumptions \\
\hline
Definitions, enlargement, empty support, equivariant images,
common product bounds, closure and transport of finite supportedness & None \\
Visible finite-set image and union bounds, swap criterion &
  \texttt{DecidableEq A} \\
Intersection of finite support bounds &
  \texttt{DecidableEq A}, \texttt{Infinite A} \\
\hline
\end{tabular}
\end{center}

All rows use the selected actions, with independent carrier universes.
Countability, carrier-wide supportedness and commuting-action assumptions
are unnecessary. The calculus therefore applies to finitely supported
elements within arbitrary permutation actions, before imposing any
condition on the carrier as a whole.

\subsection{Nominality and least support}
\label{sec:least-support}

\begin{definition}[Nominality certificate]
A selected action on $X$ is \emph{nominal} when every $x\in X$ is finitely
supported. The certificate asserts
\[
  \forall x\in X,\quad
  \exists S\in\Finset(A),\quad\operatorname{Supports}(S,x).
\]
\end{definition}

This condition is meaningful for any atom carrier. In Lean,
\texttt{Nominal A X} is a class in \texttt{Prop}, parameterized by the
already selected \texttt{MulAction (Perm A) X}. Its sole field proves finite
supportedness of every element. It contains neither an action nor a chosen
finite set, and requires neither infinitude nor decidable equality. Atom and
carrier universes remain independent; the atom parameter is explicit.

Least support concerns an individual element. Write
$\operatorname{supp}_A(x)$ for its least finite supporting set when this
exists. The next result requires finite supportedness of $x$, without a
nominality certificate for all of $X$.

\begin{proposition}[Unique least finite support]
\label{prop:least-support}
Assume $A$ is infinite and $x\in X$ is finitely supported. There is a unique
finite set $L\subseteq A$ such that
\[
  \operatorname{Supports}(L,x)
  \quad\text{and}\quad
  \forall T\in\Finset(A),\quad
  \operatorname{Supports}(T,x)\Longrightarrow L\subseteq T.
\]
Consequently, for every finite $S\subseteq A$,
\[
  \operatorname{Supports}(S,x)
  \quad\Longleftrightarrow\quad\operatorname{supp}_A(x)\subseteq S.
\]
\end{proposition}

\begin{proof}
Strict inclusion of finite sets is well founded. Since $x$ has a finite
support, choose an inclusion-minimal supporting set $L$. For any finite
support $T$, Proposition~\ref{prop:finite-support-intersection} shows that
$L\cap T$ also supports $x$. Its containment in $L$ and the minimality of
$L$ give $L\subseteq L\cap T\subseteq T$. Thus $L$ is least, rather than
merely minimal. Two least supporting sets contain each other and hence
are equal. The final equivalence follows from leastness and enlargement
of supporting bounds.
\end{proof}

Mathlib supplies the well-founded Finset order and
\path{exists_minimal_of_wellFoundedLT}; the nominal step is closure under
binary intersection. The declaration
\path{FinitelySupported.exists_least_support} combines these results.
One could instead minimize the cardinality of a supporting set: intersecting
it with another support cannot decrease its cardinality, so the intersection
equals the selected set. The inclusion-order proof uses the existing order
interface directly. It neither enumerates atoms nor assumes a theorem about
arbitrary intersections.

Let \texttt{hx} be a proof of \texttt{FinitelySupported A x}.
The noncomputable operation \texttt{hx.support} chooses the uniquely
characterized set. Its public laws
state that this set supports $x$ and is contained in every finite support;
clients need not unfold the choice. No public decidable-equality argument is
required: classical equality is local to the intersection proof. Choosing a
semantic support is not an algorithm for computing it.

Two certificates of finite supportedness, even if proved using different
finite bounds, yield the same least set. Lean's proof irrelevance identifies
the certificate arguments, and the mathematical uniqueness result identifies
any independently constructed least support. Neither explanation identifies
supports belonging to different selected actions.

For a nominal carrier, the convenient operation \texttt{support A x} uses
the class's finite-supportedness proof as its elementwise certificate.
It makes no second choice. The agreement law \texttt{support\_eq} equates
this value with \texttt{hx.support} for any supplied certificate of $x$.
Thus two nominality certificates for the same action also yield equal
supports. The atom type is an explicit argument to the carrier operation
and its laws, allowing clients to select the intended action when several
atom types occur in a context.

The carrier interface makes support available uniformly, while the
elementwise interface retains the weaker hypothesis needed for a particular
value. For instance, under the ordinary pointwise action on functions
$\mathbb N\to\mathbb N$, a constant function is supported by its value's
singleton. The identity function has no finite support: given any proposed
bound, swap two distinct atoms outside it and evaluate at one endpoint.
Consequently, this function carrier is not nominal even though its supported
elements admit the least-support construction. No function-space action is
changed to obtain this distinction.

The atom hypothesis cannot be replaced by carrier-wide nominality. Under the
canonical action on $\mathbb B$, every atom has a singleton supporting bound.
Yet the element $\mathrm{false}$ of
Example~\ref{ex:finite-atom-intersection} has no least finite support: a least
bound would lie in both singleton supports, hence be empty, whereas the empty
set does not support it. Finite supportedness and least-support existence
therefore have different atom assumptions.

\begin{proposition}[Laws of least support]
\label{prop:least-support-laws}
Assume $A$ is infinite, $x\in X$ is finitely supported, and $\pi\in\Perm(A)$.
Then
\begin{align*}
  \operatorname{supp}_A(\pi\act x)&=\pi[\operatorname{supp}_A(x)],\\
  \operatorname{supp}_A(x)=\varnothing
    &\quad\Longleftrightarrow\quad
      \forall\sigma\in\Perm(A),\ \sigma\act x=x.
\end{align*}
If $f:X\to Y$ is equivariant, then $f(x)$ is finitely supported and
\[
  \operatorname{supp}_A(f(x))\subseteq\operatorname{supp}_A(x).
\]
No nominality assumption on either whole carrier is needed.
\end{proposition}

\begin{proof}
Transport of sufficient bounds shows that
$\pi[\operatorname{supp}_A(x)]$ supports $\pi\act x$, giving one inclusion
by leastness. For the reverse inclusion, transport the least support of
$\pi\act x$ back by $\pi^{-1}$. It supports $x$, so contains
$\operatorname{supp}_A(x)$; applying $\pi$ yields the other inclusion.
Thus both directions reuse Proposition~\ref{prop:support-transport}.
For the empty characterization, the empty set is a supporting bound exactly
for elements fixed by every permutation, and it is then the least bound.
Finally, equivariance makes every support of $x$ a support of $f(x)$.
Apply this fact to $\operatorname{supp}_A(x)$ and use leastness of the image.
\end{proof}

The elementwise Lean laws take the original certificate \texttt{hx} and
derive those for the renamed element and equivariant image using
\texttt{hx.smul} and \texttt{hx.map}. A separately supplied certificate
gives the same support. The finite-set image equation uses the existing
\texttt{Pointwise} scope and requires \texttt{DecidableEq A}; the empty
characterization and image inclusion need no equality instance. None of
the laws assumes that permutation actions commute.

The carrier versions use the same elementwise laws. In particular, writing
both sides of the image bound with \texttt{support A} requires nominality of
both carriers. For an arbitrary acted-on codomain, the elementwise image
certificate still suffices. Infinitude is needed for both choices of interface;
decidable equality remains confined to the finite-set image equation.

The image inclusion is generally strict. An atom moved by a swap has
nonempty least support, while its image under a constant map into explicit
discrete data is fixed by every permutation and has empty support. Thus
equivariance preserves supporting bounds but need not preserve dependence
on every atom in the original least bound.

Least support is also distinct from strong support. If $\pi\act x=x$,
transport gives
$\pi[\operatorname{supp}_A(x)]=\operatorname{supp}_A(x)$.
This is setwise preservation, which allows atoms inside the support to be
permuted. The supporting property says that fixing those atoms pointwise
fixes $x$; it does not give the converse.
Example~\ref{ex:unordered-pair-support} makes this distinction concrete using
the exact least support of a finite atom set.
